% ===== 原文 pp.6–7 =====
% [原文 p.6]

% [环境契约 · 经 Lead 裁定] 本块从 p.5 命题 3 证明里的 itemize 中途续写：
% tr01a 在 p.5 末句 "…linear injection on each fiber." 处停下，不闭合 itemize/proof；
% 本块以 \item 开头，并在 p.6 该 bullet 段末用 \end{itemize} + \end{proof} 收尾
% （命题 3 证明的 □ 出现在 p.6 的 [单射性.] 段末）。

\item {[单射性.]} 若 $E\cong f_0^*E_n(\C^\infty), f_1^*E_n(\C^\infty)$，设 $g_0,g_1: E\to\C^\infty$ 为在每条纤维 $p^{-1}(x)$ 上都是线性单射的映射，正如证明开头所讨论的。我们要证明 $g_0\simeq g_1$，且同伦 $g_t$ 在每条纤维上都是线性单射，从而 $f_0\simeq f_1$，同伦由 $f_t(x) = g_t(p^{-1}(x))$ 给出。令 $A_t, B_t: \C^\infty\to\C^\infty$ 为同伦
\[
\begin{gathered}
A_t(z_i) = (1-t)(z_1,z_2,z_3,\dots) + t(z_1,0,z_2,0,z_3,0,\dots),\\
B_t(z_i) = (1-t)(z_1,z_2,z_3,\dots) + t(0,z_1,0,z_2,0,z_3,0,\dots).
\end{gathered}
\]
我们有 $g_0\simeq A_1\circ g_0 =: G_0$（把 $g_0$ 放入奇坐标）、$g_1\simeq B_1\circ g_1 =: G_1$（把 $g_1$ 放入偶坐标），并且 $G_0\simeq G_1$，同伦为 $(1-t)G_0 + tG_1$（全程经过的映射在每条纤维上都是线性单射）。
\end{itemize}
\end{proof}

% [Tutor Note: p.6 作者跳过了「$(1-t)G_0+tG_1$ 仍逐纤维单射」的验证，补出如下。$A_1$ 把 $(z_1,z_2,\dots)$ 送到 $(z_1,0,z_2,0,\dots)$，即把 $z_i$ 摆到第 $2i-1$ 个（奇）坐标；$B_1$ 把 $z_i$ 摆到第 $2i$ 个（偶）坐标。于是 $G_0=A_1\circ g_0$ 的像落在奇坐标子空间 $V_{\text{奇}}$ 内、$G_1=B_1\circ g_1$ 的像落在偶坐标子空间 $V_{\text{偶}}$ 内，而 $\C^\infty = V_{\text{奇}}\oplus V_{\text{偶}}$。若 $(1-t)G_0(v)+tG_1(v)=0$，按此直和分解比较两个分量，得 $(1-t)G_0(v)=0$ 且 $tG_1(v)=0$：当 $0<t<1$ 时由 $G_0$、$G_1$ 都是单射得 $v=0$；当 $t=0$ 或 $t=1$ 时分别由 $G_0$ 或 $G_1$ 单射得 $v=0$。]

这里还有 $K(X)$（$X$ 紧）的另一种分裂，它在 $X$ 连通时与 $K(X) = \widetilde{K}(X)\oplus\Z$ 一致。函数 $d: \Vect_\C(X)\to [X,\N]$，$d_E(x) = \dim p^{-1}(x)$，可延拓为 $\bar d: K(X)\to [X,\Z]$。记 $\bar d$ 的核为 $\widehat K(X)$。我们有 $K(X) = \widehat K(X)\oplus [X,\Z]$（给定 $f: X\to\Z$，$X$ 分解为有限并 $X = \bigcup_n f^{-1}(n)$，我们在 $f^{-1}(n)$ 上取维数为 $n$ 的平凡虚拟丛，即得一个分裂）。

记 $\Vect^s_\C(X)$ 为 $\Vect^k_\C(X)$ 在包含 $[E]\mapsto [E\oplus\epsilon]$ 之下的有向极限，其中 $\epsilon$ 是 $X$ 上的一维平凡丛。

% [Tutor Note: p.6 本段同时出现两个相近记号，勿混：$\widetilde{K}(X)$ 是 p.5 定义的约化 $K$ 理论（$X$ 连通时 $K(X)=\widetilde{K}(X)\oplus\Z$），而 $\widehat K(X):=\ker\bar d$ 是本段新引入的群；「另一种分裂」与旧分裂一致，正是指两者在 $X$ 连通时同构。核对方式：$\bar d$ 是群同态（$d$ 逐纤维取维数，对直和可加），且 $\bar d$ 在 $\ker\bar d$ 上取 $0$；分裂的截面正是 $f\mapsto$ 在 $f^{-1}(n)$ 上取 $n$ 维平凡虚拟丛，此时 $\bar d$ 复合该截面等于 $f$，故 $K(X)=\ker\bar d\oplus[X,\Z]$。]

\begin{lem}{1} 若 $X$ 是紧 Hausdorff 空间，$E\xrightarrow{p} X$ 是复向量丛，则存在另一个丛 $E'\xrightarrow{p'} X$ 使得 $E\oplus E'$ 是平凡的。\end{lem}

\begin{proof}
设 $X = \cup U_i$ 是有限覆盖，使得 $p^{-1}(U_i)\cong U_i\times\C^{n_i}$，并设 $\{\phi_i\}_i$ 是从属于 $\{U_i\}_i$ 的单位分解。定义 $g_i: E\to\C^{n_i}$ 为 $g(x,v) = \phi_i(x)v$，$(x,v)\in U_i\times\C^{n_i}$（于是 $g_i$ 在 $p^{-1}(U_i)$ 之外为零），并把它们用作函数 $g: E\to\bigoplus_i \C^{n_i}$ 的坐标。这个 $g$ 在每条纤维上是线性单射，而映射 $G: E\to X\times\bigoplus_i\C^{n_i}$，$G(e) = (p(e), g(e))$ 把 $E$ 实现为平凡丛的子丛。逐纤维取正交补即得所需的 $E'$。
\end{proof}

% [Tutor Note: p.6 引理 1 的证明省略了两处论证。其一，$g$ 逐纤维单射：固定 $x\in X$，由单位分解有某个 $i$ 使 $\phi_i(x)\ne 0$，而 $g$ 的第 $i$ 个分量在纤维 $p^{-1}(x)\cong\C^{n_i}$ 上就是 $\phi_i(x)$ 乘该点的坐标，是非零标量倍，故为单射，从而 $g$ 在整条纤维上单射。其二，正交补仍是丛：复向量丛的纤维带 Hermite 内积，逐纤维取 $E$ 的正交补再光滑化即得 $E'$，且 $E\oplus E'$ 平凡。另：PDF 文字层把本页的 Hausdorff 抽成了乱码，正文按 200 dpi 页图作 Hausdorff。]

最后，我们有

\begin{thm}{1}（[AGP] 9.4.6-9.4.9）若 $X$ 是紧的，则
\[ \widehat K(X)\cong \Vect^s_\C(X)\cong [X, BU], \]
由此可得
\[ K(X)\cong [X, BU\times\Z]. \]
\end{thm}

\begin{proof}
对每个 $k\ge 0$ 有映射
\[ \phi_k: \Vect^k_\C(X)\to\widehat K(X),\quad \phi_k([E]) = [E] - [\epsilon^k], \]
它与定义有向极限 $\Vect^s_\C(X)$ 的那些映射交换，因此给出映射 $\phi: \Vect^s_\C(X)\to\widehat K(X)$。我们要证明 $\phi$ 是半群同构（从而是 Abel 群同构）。对 $[E_1]-[E_2]\in\widehat K(X)$，由前一条引理，存在 $E'_i$ 使得 $E_i\oplus E'_i\cong \epsilon^n$ 是平凡的。我们有
\[ [E_1\oplus E'_2] - [\epsilon^n] = [E_1]-[E_2]\in\widehat K(X), \]
% [原文 p.7]
于是 $d([E_1\oplus E'_2]) = d([\epsilon^n]) = n$ 是常数，且 $\phi_n([E_1\oplus E'_2]) = [E_1]-[E_2]$。故 $\phi$ 是满射。

为证 $\phi$ 的单射性，若 $[E]-[\epsilon^k] = [E']-[\epsilon^l]$ 在 $\widehat K(X)$ 中（注意有向极限的每个元素都在某个 $\Vect^k_\C(X)$ 中有代表元），则存在 $n$ 使得 $E\oplus \epsilon^{l+n}\cong E'\oplus \epsilon^{k+n}$。[我们知道存在某个丛 $F$ 使得 $E\oplus \epsilon^l\oplus F\cong E'\oplus \epsilon^k\oplus F$，再添加某个满足 $F'\oplus F\cong \epsilon^n$ 的 $F'$，即得所求。] 于是 $E$ 与 $E'$ 在 $\Vect^s_\C(X)$ 中代表同一个元素。

最后，我们有
\[ \widehat K(X)\cong \Vect^s_\C(X) = \operatorname{colim}_k \Vect^k_\C(X)\cong \operatorname{colim}_k [X, BU_k]\cong [X, BU], \]
最后一个同构是因为 $X$ 紧，并且有向极限之间的同构与拉回交换（包含 $i_k: BU_k\to BU_{k+1}$ 满足 $i_k^*(E_k(\C^\infty)) = E_k(\C^\infty)\oplus \epsilon$，而包含 $\Vect^k_\C(X)\to\Vect^{k+1}_\C(X)$ 把 $[E]$ 送到 $[E\oplus\epsilon]$）。
\end{proof}

% [Tutor Note: p.7 满射性一步的两处细节。其一，原文此处写 $d([E_1\oplus E'_2])$，但按 p.6 的定义 $d$ 只定义在 $\Vect_\C(X)$ 上，$K(X)$ 上的延拓才是 $\bar d$；此处记号省了横线，应读作 $\bar d$（其值 $n$ 与 $\phi_n$ 的定义域 $\Vect^n_\C(X)$ 相符，可自洽核对）。其二，「半群同构，从而是 Abel 群同构」的理由：$\Vect^s_\C(X)$ 在 $\oplus$ 下只是交换幺半群，$\widehat K(X)$ 是 Abel 群；半群同构是保持 $\oplus$ 的双射，故源经 $\phi^{-1}$ 从靶获得逆元而成为群，$\phi$ 便是群同构。]

% [Tutor Note: p.7 单射性最后一步省略了有向极限的等价关系。由 $E\oplus\epsilon^{l+n}\cong E'\oplus\epsilon^{k+n}$ 得 $[E]=[E']$ 于 $\Vect^s_\C(X)$，是因为在 $\Vect^s_\C(X)$ 中 $[E]=[E\oplus\epsilon^{l+n}]$（往左端加了 $l+n$ 个 $\epsilon$）、$[E']=[E'\oplus\epsilon^{k+n}]$（往右端加了 $k+n$ 个 $\epsilon$），两端由同一个同构相连。方括号中 $F,F'$ 的存在性即引理 1。]

% [Tutor Note: p.7 「$\operatorname{colim}_k[X,BU_k]\cong[X,BU]$ 用到 $X$ 紧」的理由：$BU=\bigcup_k BU_k$ 是 CW 复形，紧空间到它的映射的像紧，而紧子集只与有限多个胞腔相交，故像落在某个 $BU_k$ 之中，于是该同构成立。拉回相容性指 $i_k^*(E_k(\C^\infty))=E_k(\C^\infty)\oplus\epsilon$ 与 $\Vect$ 侧 $[E]\mapsto[E\oplus\epsilon]$ 的包含相对应，这是有向极限（colim）两侧能对齐的原因。]

\section{Bott 周期性与作为广义上同调理论的 $K$ 理论}

下面一节所参考的文献是 [AGP] 第 9 章与第 12 章。

\begin{defn}{5}（[AGP] 12.1.4）可加约化上同调理论 $\widetilde{E}^*$ 是从带基点拓扑空间范畴到 Abel 群范畴的一列反变函子 $\{\widetilde{E}^n\}_{n\in\Z}$，具有以下性质：
\begin{itemize}[leftmargin=1.6em,itemsep=2pt]
\item （悬挂同构）$\widetilde{E}^n(\Sigma X)\cong \widetilde{E}^{n-1}(X)$，关于 $X$ 自然，
\item （同伦不变性）若 $f\simeq g: (X,x_0)\to (Y,y_0)$，则 $f^* = g^*: \widetilde{E}^*(X)\to \widetilde{E}^*(Y)$，
\item （正合性）对每个空间对 $(X,A)$ 存在正合列
\[ \widetilde{E}^n(X\cup CA)\to \widetilde{E}^n(X)\to \widetilde{E}^n(A) \]
（映射由包含诱导，$CA$ 是 $A$ 上的约化锥），
\item （可加性）$\widetilde{E}^n(\bigvee_\alpha X_\alpha)\cong \prod_\alpha \widetilde{E}^n(X_\alpha)$，由各包含映射诱导。
\end{itemize}
\end{defn}

% [Tutor Note: p.7 定义 5 第二条原文写 $f^*=g^*:\widetilde{E}^*(X)\to\widetilde{E}^*(Y)$，方向写反了：上同调是反变的，$f,g:(X,x_0)\to(Y,y_0)$ 诱导的是 $f^*,g^*:\widetilde{E}^*(Y)\to\widetilde{E}^*(X)$。正文照原文译，此处指出。]
% [Tutor Note: p.7 术语首现：additive reduced cohomology theory = 可加约化上同调理论；$\Omega$-prespectrum = $\Omega$-预谱；weak homotopy equivalence = 弱同伦等价；reduced cone = 约化锥；representability = 可表示性。定义 5 的四条性质是 p.8 起 Bott 周期性论证的框架：悬挂同构给出负指标的 $\widetilde{K}^{-n}$，正合性给出 Mayer–Vietoris 型长正合列。注意第三条中 $X\cup CA$ 是把 $A$ 上的（约化）锥贴到 $X$ 上。]

上同调理论的可表示性由下述对象承载（证明从略）。

\begin{defn}{6}$\Omega$-预谱是一列带基点空间 $\{E_n\}_{n\in\Z}$ 以及弱同伦等价 $\epsilon_n: E_n\to \Omega E_{n+1}$。\end{defn}

% [Tutor Note: p.7 原文定义 6 把 prespectrum 拼成 "presepectrum"（pres 后多插了一个 e），经 1400 dpi 裁图逐字复核确认（非连字误读），正文按其数学含义译作「$\Omega$-预谱」，拼写可疑处记入译注汇总。另注：$\Omega$ 是环路空间函子，$\epsilon_n: E_n\to\Omega E_{n+1}$ 是弱同伦等价，故预谱给出的 $\widetilde{E}^n=[X,E_n]$ 自带悬挂同构（因为 $[X,\Omega E_{n+1}]\cong[\Sigma X,E_{n+1}]$）。]

下面定理的困难方向来自 Brown 可表示性定理，[AGP] 定理 12.2.22。

\begin{thm}{2}（[AGP] 12.3.2, 12.3.3）一个 $\Omega$-预谱通过
\[ \widetilde{E}^n = [X, E_n]_* \]
决定了带基点拓扑空间范畴上的一个可加约化上同调理论 $\widetilde{E}^*$。反之，带基点 CW 复形范畴上的一个可加约化上同调理论 $\widetilde{E}^*$ 决定了 $\Omega$-预谱 $\{E_n\}_{n\in\Z}$，使得
\[ \widetilde{E}^n = [X, E_n]_*. \]
\end{thm}

% [Tutor Note: p.7 定理 2 中两处原文拼写笔误（均经 1100–1500 dpi 裁图逐字复核，已排除 fi/fl 连字误读，正文照原意译、不改原文）：(i) 首句 "cohmol-ogy"（原文在行末断作 cohmol-、下行接 ogy）应为 "cohomology"；紧接的下一段同一短语又写对了 "cohomology"，两处互证，可作判断依据。(ii) 后半 "additve" 应为 "additive"。另有两处非拼写问题：一是 "Representability of cohomology theories are mediated" 主谓不一致，应为 "is mediated"；二是 "proofs definitely omitted" 并非笔误（文字层作 defi-/nitely 只是 fi 连字，印刷确为 definitely），而是作者口语化的戏言，译作「证明从略」。]
